% TOP_PROBLEM: {"id": 2911, "title": "Kirby Problem 4.35", "category_id": 7, "category": {"id": 7, "name": "topology", "display_name": "Topology", "description": "Properties preserved under continuous deformations.", "slug": "topology", "order_index": 7, "created_at": "2026-07-31T15:26:25.671Z"}, "status": "open", "problem_number": "KP-4.35", "set_id": 11, "statusQualification": "Corrects the Laurent-ring transcription and makes the deck-generator convention explicit."}
% FIRST_POSTED: 2026-10-01
% ENTRY_KIND: statement_only
% UnsolvedMath ID 2911; source catalogue code KP-4.35
% !TeX program = lualatex
% Definition and sources only; this file is not an LLM solution attempt.
\documentclass[11pt,a4paper]{article}
\usepackage[margin=25mm]{geometry}
\usepackage{fontspec}
\setmainfont{lmroman10-regular.otf}[BoldFont=lmroman10-bold.otf,
 ItalicFont=lmroman10-italic.otf,BoldItalicFont=lmroman10-bolditalic.otf]
\usepackage{amsmath,amssymb,mathtools}
\usepackage{unicode-math}
\setmathfont{latinmodern-math.otf}
\AtBeginDocument{\let\setminus\smallsetminus}
\usepackage{xurl,hyperref}
\hypersetup{colorlinks=true,urlcolor=blue}
\setcounter{secnumdepth}{0}
\setlength{\parindent}{0pt}
\setlength{\parskip}{5pt}
\setlength{\emergencystretch}{3em}
\begin{document}
\section*{Kirby Problem 4.35}\label{2911}
{\small UnsolvedMath ID 2911 · KP-4.35 · Topology · Kirby's Problems in Low-Dimensional Topology}
\subsection{Definitions and mathematical statement}
A two-knot is an embedding $k:S^2\hookrightarrow S^4$, either smooth or locally flat topological; address the two categories separately. Its exterior is $E_k=S^4\setminus\operatorname{int}\nu k(S^2)$, where $\nu k(S^2)\cong S^2\times D^2$. The exterior and the open complement have the same homotopy type. Its knot group is $G_k=\pi_1(E_k)$.

(a) Give intrinsic algebraic necessary and sufficient conditions on a group $G$ for $G\cong G_k$ for a two-knot in the standard $S^4$. This asks for a structural characterization, beyond restating existence of a knot diagram or a specially realized presentation.

An orientation of the knot and ambient sphere identifies $H_1(E_k;\mathbb Z)\cong\mathbb Z$, generated by a positive meridian. Let $\widetilde E_k$ be the infinite cyclic cover corresponding to the kernel of this abelianization. A generator $t$ of its deck transformations makes
\[
\mathcal A_k=H_1(\widetilde E_k;\mathbb Z)
\]
a module over $\Lambda=\mathbb Z[t,t^{-1}]$.

(b) Classify the $\Lambda$-modules isomorphic to $\mathcal A_k$ for such knots, again distinguishing smooth and locally flat realizability. Reversing the meridian convention replaces the $t$ action by its inverse; specify whether this identification is allowed in a classification.

Integral module structure is required, rather than only its rationalization or its order polynomial. Likewise, realization in a homotopy four-sphere alone does not immediately give smooth realization in the standard sphere.
\subsection{Short English statement}
Characterize two-knot groups and integral Alexander modules in the standard four-sphere.
\subsection{Sources}
\begin{itemize}
\item[S1] ULAM.AI and the UnsolvedMath Contributors, supplied problems.json, record 2911 (KP-4.35), Kirby's Problems in Low-Dimensional Topology. Catalogue identity and source transcription; CC BY 4.0.
\url{https://huggingface.co/datasets/ulamai/UnsolvedMath}.
\item[S2] K3: A New Problem List in Low-Dimensional Topology, Problem 4.35. Expanded formulation of the supplied catalogue statement.
\url{https://math.berkeley.edu/sites/default/files/surv-295-ruberman-watermarked-author-pdf.pdf}.
\end{itemize}
\end{document}
